\scanpage{73}
\mspage{4}
(9) de passage de $f$ à $g$, sous une forme considérablement plus simple.

Soit $T$ un espace topologique \deleted{simplement} $1$-connexe, et
\[
(13)\qquad p:T\longrightarrow X
\]
une application continue \inserted{(le cas \unsure{injectif} sera le plus intéressant)}. Soit $T_K=p^{-1}(X_K)$.

Considérons \deleted{\unsure{$\Pi_1(T_K),\ \Pi_1$}}
\deleted{\unsure{Soient} $x,y\in T_K$}
\[
\unsure{(14)}\qquad\deleted{\unsure{\Pi_1(X)}}\ \Pi_1(X;T)=\pi
\]
alors pour tout $x\in T$ on a un isom canonique
\[
(15)\qquad\varphi_x:\pi\xrightarrow{\sim}\Pi_1(X,p(x))
\]
\inserted{de sorte que les $\Pi_1(X,p(x))$ sont munis d'un syst. transitif d'iso}
\margin{On désigne \unsure{ainsi|aussi} par $\varphi_x$ le prolongement aux complétés profinis\ldots}
et pour $x,y\in T$, l'unique \deleted{élément} \inserted{flèche} \deleted{de Hom} $x:y$ de $\Pi_1(T)$ donne \inserted{par $p$} une flèche canonique
\[
\unsure{(16)}\qquad l_{y,x}:p(x)\mathbin{-}p(y)\qquad\text{dans }\Pi_1(X)
\]
(NB cette flèche dépend du choix de $x,y$ au dessus de $px,py$ --- mais naturellement $p$ sera injective).

Moyennant les iso (15), on peut considérer $f_{l_{y,x}}(u)$ et $g_{l_{y,x}}(u)$ (\inserted{du moins} si $x,y\in T_K$ i.e. $p(x),p(y)\in X_K$) comme des éléments du $=$ groupe $\pi$ \unsure{(14)}, \deleted{soit} et en fait ces deux éléments se réduisent à un seul et $=$ élément, soit
\[
\begin{array}{c}
\unsure{(17)}\qquad g_{y,x}^{T/X}(u)=g_{y,x}^{p}(u)=\varphi_y^{-1}\bigl(g_{l_{y,x}}(u)\bigr)=\varphi_x^{-1}\bigl(f_{l_{y,x}}(u)\bigr)\\
\rotatebox{90}{$\in$}\\
\pi=\Pi_1(X;T)
\end{array}
\]
\margin{72}
